\chapter{Matemática}

\section{Operaciones basicas}

\subsection{Obtener $\pi$}
\begin{lstlisting}[language=C++, caption={Obtener $\pi$}]
	const double PI = 3.141592653589793;  //FORMA 1
	double pi = M_PI; 	 //FORMA 2
	
\end{lstlisting}

\subsection{Obtener la raiz cuadrada}
Para calcular la raíz cuadrada de un número en C++, se utiliza la función \texttt{sqrt()}.

\bigskip

\textbf{Ejemplo:}

\begin{lstlisting}[language=C++]
	double x = 25.0;
	double raiz = sqrt(x);
\end{lstlisting}

\bigskip

\textbf{Notas importantes:}
\begin{itemize}
	\item El resultado puede guardarse en un \texttt{double} para conservar la precisión decimal.
	\item Si se guarda en un tipo entero, la parte decimal se truncará (se corta sin redondeo).
\end{itemize}

Complejidad: \( O(\sqrt{n}) \)
\begin{lstlisting}[language=C++, caption={Chequear si un número es primo}]
	bool esPrimo(int n) {
		if (n <= 1) return false;
		if (n == 2) return true;
		if (n % 2 == 0) return false;
		for (int i = 3; i * i <= n; i += 2)
		if (n % i == 0)
		return false;
		return true;
	}
\end{lstlisting}

\section{Reglas de Divisibilidad}

Para determinar si un número es divisible por otro sin necesidad de hacer la división, se pueden usar las siguientes reglas de divisibilidad:

\subsection{Reglas de Divisibilidad Básicas}

\begin{enumerate}
	\item \textbf{Divisible por 1}: Todos los números son divisibles por 1.
	\item \textbf{Divisible por 2}: Un número es divisible por 2 si su último dígito es par (0, 2, 4, 6, 8).
	\item \textbf{Divisible por 3}: Un número es divisible por 3 si la suma de sus dígitos es divisible por 3.
	\item \textbf{Divisible por 4}: Un número es divisible por 4 si el número formado por sus últimos dos dígitos es divisible por 4.
	\item \textbf{Divisible por 5}: Un número es divisible por 5 si su último dígito es 0 o 5.
	\item \textbf{Divisible por 6}: Un número es divisible por 6 si es divisible por 2 y por 3.
	\item \textbf{Divisible por 7}: Hay varias reglas:
	\begin{itemize}
		\item \textbf{Regla 1}: Toma el último dígito, multiplícalo por 2 y réstalo del resto del número. Si el resultado es divisible por 7, entonces el número original también lo es.
		\item \textbf{Regla 2}: Divide el número en bloques de 3 dígitos (de derecha a izquierda) y calcula la suma alternante de estos bloques. Si el resultado es divisible por 7, entonces el número original también lo es.
	\end{itemize}
	\item \textbf{Divisible por 8}: Un número es divisible por 8 si el número formado por sus últimos tres dígitos es divisible por 8.
	\item \textbf{Divisible por 9}: Un número es divisible por 9 si la suma de sus dígitos es divisible por 9.
	\item \textbf{Divisible por 10}: Un número es divisible por 10 si su último dígito es 0.
\end{enumerate}

\subsection{Reglas de Divisibilidad Menos Comunes}

\begin{enumerate}
	\item \textbf{Divisible por 11}:
	\begin{itemize}
		\item \textbf{Regla 1}: Suma los dígitos en posiciones impares y resta la suma de los dígitos en posiciones pares. Si el resultado es divisible por 11, entonces el número original también lo es.
		\item \textbf{Regla 2}: Alternativamente, puedes usar la suma alternante de los dígitos.
	\end{itemize}
	\item \textbf{Divisible por 12}: Un número es divisible por 12 si es divisible por 3 y por 4.
	\item \textbf{Divisible por 13}:
	\begin{itemize}
		\item \textbf{Regla 1}: Toma el último dígito, multiplícalo por 9 y réstalo del resto del número. Si el resultado es divisible por 13, entonces el número original también lo es.
		\item \textbf{Regla 2}: Divide el número en bloques de 3 dígitos y calcula la suma alternante de estos bloques. Si el resultado es divisible por 13, entonces el número original también lo es.
	\end{itemize}
	\item \textbf{Divisible por 17}:
	\begin{itemize}
		\item \textbf{Regla 1}: Toma el último dígito, multiplícalo por 5 y réstalo del resto del número. Si el resultado es divisible por 17, entonces el número original también lo es.
	\end{itemize}
	\item \textbf{Divisible por 19}:
	\begin{itemize}
		\item \textbf{Regla 1}: Toma el último dígito, multiplícalo por 2 y súmalo al resto del número. Si el resultado es divisible por 19, entonces el número original también lo es.
	\end{itemize}
\end{enumerate}

\subsection{Reglas de Divisibilidad para Números Primos Mayores}

Para números primos mayores, como 23, 29, etc., las reglas pueden volverse más complejas y a menudo implican métodos similares a los usados para 7, 13, 17 y 19, pero con diferentes multiplicadores.

\subsection{Métodos Generales para Verificar Divisibilidad}

\begin{enumerate}
	\item \textbf{División Directa}: La forma más directa de verificar la divisibilidad es realizar la división y ver si el residuo es cero.
	\item \textbf{Uso de Factores}: Si un número es divisible por dos o más números coprimos, entonces es divisible por su producto. Por ejemplo, si un número es divisible por 3 y por 5, entonces es divisible por 15.
\end{enumerate}

\subsection{Generación de divisores}
Para encontrar todos los divisores de un número $k$:
\begin{lstlisting}[language=C++, caption={Generar los divisores de un numero}]
	for (ll i = 1; i * i <= k; i++) {
		if (k % i == 0) {
			divs.push_back(i);
			if (i != k / i) divs.push_back(k / i);
		}
	}
\end{lstlisting}

\section{Números Primos}

\subsection{Chequear si un número es primo}

Complejidad: \( O(\sqrt{n}) \)
\begin{lstlisting}[language=C++, caption={Chequear si un número es primo}]
	bool esPrimo(int n) {
		if (n <= 1) return false;
		if (n == 2) return true;
		if (n % 2 == 0) return false;
		for (int i = 3; i * i <= n; i += 2)
		if (n % i == 0)
		return false;
		return true;
	}
\end{lstlisting}

\subsection{Criba de Eratóstenes para factorizar números}

\textbf{Uso} \\
Guarda para cada número su \textbf{menor factor primo} (SPF, \textit{smallest prime factor}).
Con eso, puedes factorizar cualquier número \(x \le n\) muy rápido.

\begin{lstlisting}[language=C++, caption={Criba SPF + factorización en primos}]

vector<int> spf; // spf[x] = menor factor primo de x

void build_spf(int n) {
	spf.assign(n + 1, 0);
	for (int i = 2; i <= n; i++) {
		if (spf[i] == 0) {
			spf[i] = i; // i es primo
			if (1LL * i * i <= n) {
				for (long long j = 1LL * i * i; j <= n; j += i) {
					if (spf[j] == 0) spf[j] = i;
				}
			}
		}
	}
}

vector<pair<int,int>> factorize(int x) {
	vector<pair<int,int>> fac;
	while (x > 1) {
		int p = spf[x], e = 0;
		while (x % p == 0) {
			x /= p;
			e++;
		}
		fac.push_back({p, e});
	}
	return fac;
}

int main() {
	int n = 100;         // limite maximo que vas a factorizar
	build_spf(n);

	int x = 6;
	auto f = factorize(x);

	// salida: 6 = 2^1 * 3^1
	cout << x << " = ";
	for (int i = 0; i < (int)f.size(); i++) {
		cout << f[i].first << "^" << f[i].second;
		if (i + 1 < (int)f.size()) cout << " * ";
	}
}
\end{lstlisting}

\textbf{Complejidad:}
\begin{itemize}
	\item Preprocesamiento de la criba SPF hasta \(n\): \(O(n \log\log n)\).
	\item Memoria: \(O(n)\).
	\item Factorizar un número \(x\): \(O(k)\), donde \(k\) es la cantidad de factores primos (con multiplicidad). En práctica: muy rápido.
\end{itemize}

\section{Permutaciones}

\subsection{Computar las permutaciones}

Sea un arreglo \(A = [a_1, a_2, \dots, a_n]\). Queremos generar todas las permutaciones de sus elementos.

\begin{lstlisting}[language=C++, caption={Chequear si un número es primo}]
	#include <bits/stdc++.h>
	using namespace std;
	
	int main() {
		vector<int> A = {1, 2, 3};
		sort(A.begin(), A.end()); // ordenar antes de usar next_permutation
		
		do {
			for (int x : A) cout << x << " ";
			cout << "\n";
		} while (next_permutation(A.begin(), A.end()));
		
		return 0;
	}
\end{lstlisting}

\section{Geometría}

\subsection{Distancia entre dos puntos}
Dados dos puntos \((x_1, y_1)\) y \((x_2, y_2)\), la distancia euclidiana es:

\[
d = \sqrt{(x_1-x_2)^2 + (y_1-y_2)^2}
\]

En programación competitiva, muchas veces conviene usar la distancia al cuadrado (para evitar \texttt{sqrt}):

\[
d^2 = (x_1-x_2)^2 + (y_1-y_2)^2
\]

\begin{lstlisting}[language=C++, caption={Distancia al cuadrado entre dos puntos}]
	ll dist2 = (x1 - x2) * (x1 - x2) + (y1 - y2) * (y1 - y2);
\end{lstlisting}

\subsection{Triángulos}
Un triángulo queda determinado si conocés:
\begin{itemize}
	\item 3 lados (SSS)
	\item 2 lados + 1 ángulo (SAS)
	\item 2 ángulos + 1 lado (ASA / AAS)
\end{itemize}

\subsubsection{Ley de Cosenos}
Se usa cuando tenés lados y querés ángulos, o viceversa:

\[
c^2 = a^2 + b^2 - 2ab\cos(C)
\]

\textbf{Sirve para:}
\begin{itemize}
	\item SSS \(\to\) sacar ángulos
	\item SAS \(\to\) sacar el lado faltante
\end{itemize}

\textbf{Ejemplo (CP):}
\begin{lstlisting}[language=C++, caption={Ley de Cosenos: obtener un ángulo}]
	double C = acos((a * a + b * b - c * c) / (2.0 * a * b));
\end{lstlisting}

\subsubsection{Ley de Senos}
Cuando tenés ángulos y lados mezclados:

\[
\frac{\sin(A)}{a} = \frac{\sin(B)}{b} = \frac{\sin(C)}{c}
\]

\textbf{Sirve para:}
\begin{itemize}
	\item ASA / AAS \(\to\) sacar lados
	\item completar triángulos
\end{itemize}

\textbf{Ejemplo:}
\begin{lstlisting}[language=C++, caption={Ley de Senos: obtener un lado}]
	double b = a * sin(B) / sin(A);
\end{lstlisting}


